\documentclass[12pt]{article}
\usepackage[utf8]{inputenc}
\usepackage[T1]{fontenc}
\usepackage[margin=1in]{geometry}
\usepackage{setspace}
\usepackage{xcolor}
\usepackage[colorlinks=true, urlcolor=blue!60!black, linkcolor=blue!60!black]{hyperref}
\setlength{\parindent}{2em}
\setlength{\parskip}{1.25em}
\renewcommand{\baselinestretch}{1.0}
\begin{document}
\begin{center}
{\LARGE \textit{United States v. Cruikshank}, 92 U.S. 542 (1876)\par}
\vspace{0.8em}
{\large Case Brief\par}
\end{center}
\vspace{1em}

\section{Memory Jogger}

Overturning the federal conspiracy convictions arising from the Colfax Massacre, the Supreme Court held that the First and Second Amendments bind only Congress and that the Fourteenth Amendment bars only state---not private---action, severely curbing federal authority to prosecute racial terror during Reconstruction.

\section{Detailed Case Facts}

The prosecution in \textit{United States v. Cruikshank} grew out of the Colfax Massacre, the deadliest single episode of racial violence during Reconstruction (\href{https://calendar.eji.org/racial-injustice/apr/13}{Equal Justice Initiative}; \href{https://civics.supremecourthistory.org/article/united-states-v-cruikshank/}{Supreme Court Historical Society}). Following Louisiana's bitterly disputed 1872 gubernatorial election, rival Republican and Democratic ``Fusionist'' slates claimed authority over state and local offices, including the sheriff's office and judgeship in Grant Parish (\href{https://storyoftheweek.loa.org/2023/04/the-colfax-massacre-trial.html}{Library of America}). Fearing an armed Democratic takeover of the parish government, Black Republican freedmen assembled to guard the Grant Parish courthouse in Colfax, Louisiana. On Easter Sunday, April 13, 1873, a heavily armed white mob---including former Confederate soldiers and white supremacist paramilitaries---attacked the courthouse with rifles and a cannon, overwhelmed the defenders, and murdered scores of Black men, including dozens executed after surrendering (\href{https://calendar.eji.org/racial-injustice/apr/13}{Equal Justice Initiative}; \href{https://storyoftheweek.loa.org/2023/04/the-colfax-massacre-trial.html}{Library of America}). Historical estimates of the Black citizens killed range from roughly 60 to 150, while three white attackers died.

Because state authorities did not prosecute the perpetrators, federal prosecutors obtained a 32-count indictment against more than one hundred men in the U.S.~Circuit Court for the District of Louisiana under the Enforcement Act of May 31, 1870, ch.~114, 16 Stat.~140 (\href{https://supreme.justia.com/cases/federal/us/92/542/}{92 U.S. at 544--48, 560}; \href{https://www.courtlistener.com/opinion/89309/united-states-v-cruikshank/}{CourtListener}). Section~6 of the Enforcement Act made it a federal felony (punishable by a fine up to \$5,000, imprisonment up to ten years, and disqualification from federal office) for two or more persons to ``band or conspire together, or go in disguise upon the public highway, or upon the premises of another,'' with intent to violate the Act or ``to injure, oppress, threaten, or intimidate any citizen, with intent to prevent or hinder his free exercise and enjoyment of any right or privilege granted or secured to him by the Constitution or laws of the United States, or because of his having exercised the same'' (\href{https://supreme.justia.com/cases/federal/us/92/542/}{92 U.S. at 548}).

The first sixteen counts charged violations of Section~6 in two parallel series of eight counts---the first eight alleging ``banding together'' and counts nine through sixteen alleging ``conspiring''---directed against two Black U.S.~citizens, Levi Nelson and Alexander Tillman (\href{https://supreme.justia.com/cases/federal/us/92/542/}{92 U.S. at 544--45, 548}). Those paired counts alleged conspiracies to hinder Nelson and Tillman in: (1/9) peaceably assembling for a lawful purpose; (2/10) bearing arms for a lawful purpose; (3/11) enjoying life and personal liberty without deprivation absent due process; (4/12) enjoying the full and equal benefit of laws for security of person and property; (5/13) exercising their federal and state citizenship rights ``on account of'' their race and color; (6/14) voting in future Louisiana elections; (7/15) having voted in the November 4, 1872 Louisiana election; and (8/16) enjoying ``every, each, all, and singular'' rights granted by the federal Constitution and laws (\href{https://supreme.justia.com/cases/federal/us/92/542/}{92 U.S. at 544--45}). Counts seventeen through thirty-two charged murder under Section~7 committed in the course of violating Sections~5 and~6 (\href{https://supreme.justia.com/cases/federal/us/92/542/}{92 U.S. at 560}).

\section{Procedural History}

\begin{itemize}
    \item \textbf{First Circuit Court Trial:} Nine defendants were brought to trial in the U.S.~Circuit Court for the District of Louisiana before Circuit Judge William B. Woods. One defendant (Alfred Lewis) was acquitted, and the jury deadlocked on the remaining eight, resulting in a mistrial (\href{https://civics.supremecourthistory.org/article/united-states-v-cruikshank/}{Supreme Court Historical Society}; \href{https://storyoftheweek.loa.org/2023/04/the-colfax-massacre-trial.html}{Library of America}).
    \item \textbf{Second Circuit Court Trial (June 1874):} Eight defendants were retried before Judge Woods and Supreme Court Associate Justice Joseph P. Bradley riding circuit. On June 10, 1874, the jury acquitted five defendants on all counts, and convicted three---William Cruikshank, John P. Hadnot, and William Irwin---on the first sixteen conspiracy counts under Section~6, while acquitting them on the remaining sixteen murder counts (\href{https://supreme.justia.com/cases/federal/us/92/542/}{92 U.S. at 545, 560--61}; \href{https://civics.supremecourthistory.org/article/united-states-v-cruikshank/}{Supreme Court Historical Society}).
    \item \textbf{Motion in Arrest of Judgment and Certificate of Division:} After their motion for a new trial was denied, the three convicted defendants moved in arrest of judgment, arguing that the sixteen counts did not state indictable federal offenses and that Section~6 exceeded Congress's constitutional power (\href{https://supreme.justia.com/cases/federal/us/92/542/}{92 U.S. at 546}). In a widely circulated circuit opinion (\textit{United States v. Cruikshank}, 25 F.~Cas.~707 (C.C.D.~La.~1874)), presiding Justice Bradley concluded that the counts were legally insufficient and that judgment should be arrested (\href{https://www.fjc.gov/history/exhibits/circuit-court-opinions/United-States-v-Cruikshank}{Federal Judicial Center}). Judge Woods disagreed, creating a formal division of opinion certified to the U.S.~Supreme Court (\href{https://supreme.justia.com/cases/federal/us/92/542/}{92 U.S. at 546, 548}).
    \item \textbf{U.S.~Supreme Court:} The case was argued in spring 1875 (October Term 1875) and decided on March 27, 1876. The Supreme Court affirmed the Circuit Court's order arresting judgment and remanded with instructions to discharge the defendants (\href{https://supreme.justia.com/cases/federal/us/92/542/}{92 U.S. at 559}).
\end{itemize}

\section{Judicial Votes}

The Court was \textbf{unanimous (9--0) in the judgment} affirming the arrest of judgment and discharging the defendants, but divided \textbf{8--1 on the reasoning} (\href{https://supreme.justia.com/cases/federal/us/92/542/}{92 U.S. at 548, 559}; \href{https://ballotpedia.org/UNITED_STATES_v._CRUIKSHANK_et_al._(1876)}{Ballotpedia}):

\begin{itemize}
    \item \textbf{Majority Opinion (8 Justices):} Chief Justice Morrison R. Waite (author), joined by Associate Justices Noah H. Swayne, Samuel Freeman Miller, David Davis, Stephen J. Field, William Strong, Joseph P. Bradley, and Ward Hunt.
    \item \textbf{Dissenting from the Opinion / Concurring in the Judgment (1 Justice):} Associate Justice Nathan Clifford filed a solo opinion styled ``Mr.~Justice Clifford dissenting'' (\href{https://supreme.justia.com/cases/federal/us/92/542/}{92 U.S. at 559--69}). Clifford explicitly concurred that judgment had to be arrested, but solely because the indictment counts were fatally vague under Sixth Amendment and common-law criminal pleading rules, declining to join the majority's constitutional restrictions on the Fourteenth Amendment and the Enforcement Act.
\end{itemize}

\section{Holding}

\textbf{Judgment:} The Circuit Court's order arresting judgment upon the verdict is affirmed, and the case is remanded with instructions to discharge the defendants (\href{https://supreme.justia.com/cases/federal/us/92/542/}{92 U.S. at 559}).

\textbf{Rule Announced:} To sustain an indictment under Section~6 of the Enforcement Act of 1870, the government must allege with particularity a conspiracy to interfere with a right granted or secured by the federal Constitution or federal laws. Under our system of dual citizenship, the federal government may protect only attributes of national citizenship:
\begin{enumerate}
    \item The First Amendment right to assemble and the Second Amendment right to bear arms restrict only Congress, not the states or private citizens; only assembly for a federal purpose (such as petitioning Congress) is an attribute of national citizenship (\href{https://supreme.justia.com/cases/federal/us/92/542/}{92 U.S. at 551--53}).
    \item The Fourteenth Amendment's Due Process and Equal Protection Clauses prohibit \textit{state} encroachment only; they add nothing to the rights of one private citizen against another, leaving ordinary protection of life, liberty, and equality to the states (\href{https://supreme.justia.com/cases/federal/us/92/542/}{92 U.S. at 553--55}).
    \item The Fifteenth Amendment confers no general federal right to vote in state elections, only exemption from racial discrimination in voting; an indictment that omits an explicit allegation of racial motivation is substantively defective (\href{https://supreme.justia.com/cases/federal/us/92/542/}{92 U.S. at 555--56}).
    \item Counts alleging a conspiracy to deny ``every, each, all, and singular'' federal rights without specifying any particular right are fatally vague under the Sixth Amendment (\href{https://supreme.justia.com/cases/federal/us/92/542/}{92 U.S. at 557--59}).
\end{enumerate}

\textbf{Holding vs.~Dicta:} The core \textbf{holding} is that none of the sixteen counts stated an offense indictable under Section~6 of the Enforcement Act, requiring arrest of judgment. However, because the government's counsel pressed only counts 13, 14, and 16 at oral argument (\href{https://supreme.justia.com/cases/federal/us/92/542/}{92 U.S. at 561} (Clifford, J., dissenting)) and because those counts were independently invalid for pleading vagueness, the majority's sweeping constitutional pronouncements---especially its categorical non-incorporation of the First and Second Amendments and its narrow reading of federal enforcement authority under the Fourteenth Amendment---reached far beyond what was strictly necessary to resolve the case. Later Supreme Court decisions superseded \textit{Cruikshank}'s non-incorporation rulings by selectively incorporating the First Amendment right of assembly (\textit{De Jonge v. Oregon}, 299 U.S. 353 (1937)) and the Second Amendment right to keep and bear arms (\href{https://www.law.cornell.edu/supct/html/08-1521.ZO.html}{\textit{McDonald v. City of Chicago}, 561 U.S. 742 (2010)}) against the states through the Fourteenth Amendment's Due Process Clause, while preserving its core state-action limitation.

\section{Analysis of Opinions}

\subsection{Opinion of the Court (Chief Justice Waite)}

Chief Justice Waite began from the premise of dual sovereignty and dual citizenship articulated three years earlier in the \textit{Slaughter-House Cases}, 83 U.S. (16 Wall.) 36 (1873) (\href{https://supreme.justia.com/cases/federal/us/92/542/}{92 U.S. at 549--51}). Because the United States is a government of delegated powers alone, a person's rights of federal citizenship differ from his rights of state citizenship, and Congress may protect only those rights granted or secured by the federal Constitution or federal statutes.

Applying that framework count by count, Waite dismantled the indictment:
\begin{itemize}
    \item \textbf{Assembly (Counts 1 \& 9):} The general right to peaceably assemble predated the Constitution and belongs to state protection. Relying on \textit{Barron v. Baltimore}, 32 U.S. (7 Pet.) 243 (1833), Waite held that the First Amendment restricts only Congress. While assembling to petition Congress or conduct national public affairs is an attribute of national citizenship, the indictment failed to allege that the Colfax meeting had any federal purpose (\href{https://supreme.justia.com/cases/federal/us/92/542/}{92 U.S. at 551--53}).
    \item \textbf{Bearing Arms (Counts 2 \& 10):} The right to bear arms for a lawful purpose is not created by the Constitution; the Second Amendment ``means no more than that it shall not be infringed by Congress,'' leaving citizens to look to state police power for protection against private interference (\href{https://supreme.justia.com/cases/federal/us/92/542/}{92 U.S. at 553}).
    \item \textbf{Life, Liberty, Due Process, and Equal Protection (Counts 3, 4, 11 \& 12):} Conspiracy to murder or falsely imprison within a state is a state offense. The Fourteenth Amendment prohibits \textit{states} from depriving persons of life or liberty without due process or denying equal protection, ``but this adds nothing to the rights of one citizen as against another'' (\href{https://supreme.justia.com/cases/federal/us/92/542/}{92 U.S. at 554--55}). The federal government's sole obligation is to ensure that states do not deny those rights. Moreover, no claim arose under the Civil Rights Act of 1866 because counts 4 and 12 did not explicitly allege that the defendants acted on account of the victims' race.
    \item \textbf{Voting Rights (Counts 6, 7, 14 \& 15):} Citing \textit{Minor v. Happersett}, 88 U.S. (21 Wall.) 162 (1875), and \textit{United States v. Reese}, 92 U.S. 214 (1876), Waite explained that the right to vote in state elections comes from the states; the Fifteenth Amendment confers only exemption from racial discrimination in voting. Although the Court ``may suspect that race was the cause of the hostility,'' the counts failed to aver racial motivation positively, and failed to specify federal elections (\href{https://supreme.justia.com/cases/federal/us/92/542/}{92 U.S. at 555--56}).
    \item \textbf{Vagueness (Counts 5, 8, 13 \& 16):} Finally, the remaining counts charged a generic conspiracy to deprive the victims of ``every, each, all, and singular'' federal rights without specifying which right was targeted. Under the Sixth Amendment right ``to be informed of the nature and cause of the accusation,'' an indictment must ``descend to particulars'' so both the accused and the court can identify the offense (\href{https://supreme.justia.com/cases/federal/us/92/542/}{92 U.S. at 557--59}).
\end{itemize}

\subsection{Dissenting Opinion / Concurrence in the Judgment (Justice Clifford)}

Justice Clifford concurred that judgment must be arrested, ``but for reasons quite different from those given by the court'' (\href{https://supreme.justia.com/cases/federal/us/92/542/}{92 U.S. at 559}). Observing that government counsel had confined their argument to counts 13, 14, and 16, Clifford limited his analysis to those three counts and resolved the case strictly on criminal pleading grounds (\href{https://supreme.justia.com/cases/federal/us/92/542/}{92 U.S. at 561}).

Clifford emphasized that federal courts possess no common-law criminal jurisdiction and that every statutory offense must be pleaded with reasonable particularity under the Sixth Amendment. Counts 13 and 16 failed because they recited generic statutory language without identifying any specific federal right invaded or any overt act (\href{https://supreme.justia.com/cases/federal/us/92/542/}{92 U.S. at 565--66, 568--69}). Count 14 identified voting as the right, but failed to specify the purpose, time, or place of the future elections (\href{https://supreme.justia.com/cases/federal/us/92/542/}{92 U.S. at 567--68}). Crucially, unlike Waite, Clifford recognized Congress's affirmative power to enforce the Fourteenth Amendment's Privileges or Immunities Clause and declined to reach the defendants' broad constitutional challenges to the Enforcement Act once the pleading defects proved dispositive (\href{https://supreme.justia.com/cases/federal/us/92/542/}{92 U.S. at 563, 569}).

\section{Examples: Future Applications}

\begin{enumerate}
    \item \textbf{Same-Side (Outside Federal Fourteenth Amendment Enforcement):} \textit{Private Mob Assault Without State Involvement or Federal Nexus.} A group of private individuals assaults a political rival in a dispute over a municipal zoning ordinance, with no participation by state officials, no racial animus, and no connection to federal elections or federal institutions. Under \textit{Cruikshank}'s state-action doctrine, the Fourteenth Amendment does not authorize federal criminal prosecution of purely private violence; remedy lies in state criminal law.
    \item \textbf{Same-Side (Invalid Indictment for Vagueness):} \textit{Generic Federal Civil-Rights Indictment.} A federal grand jury indicts two defendants under 18 U.S.C.~\S~241 (the modern successor to Section~6) simply by tracking the statute and alleging they conspired ``to deprive Victim of all rights secured by the Constitution and laws of the United States,'' without naming the specific right or factual circumstances. Under \textit{Cruikshank}'s Sixth Amendment pleading rule (endorsed by both Waite and Clifford), the indictment must be dismissed as fatally vague.
    \item \textbf{Opposite-Side (Valid Federal Prosecution):} \textit{Private Conspiracy to Obstruct a Federal Petition or Congressional Election.} Private paramilitaries violently disperse citizens gathering to sign a petition to Congress or intimidate Black voters from casting ballots in a U.S.~House election because of their race. Even under \textit{Cruikshank} itself, assembling to petition Congress and exercising federal voting rights free from racial discrimination are attributes of national citizenship that Congress may protect against private conspiracies (\href{https://supreme.justia.com/cases/federal/us/92/542/}{92 U.S. at 552--53, 555--56}).
    \item \textbf{Opposite-Side (Valid Federal Enforcement via State Action):} \textit{State Officials Conspiring with Private Citizens.} A county sheriff and deputies conspire with private Klansmen to release civil-rights activists from jail so the mob can murder them (as later occurred in \textit{United States v. Price}, 383 U.S. 787 (1966), and \href{https://www.law.cornell.edu/supremecourt/text/383/745}{\textit{United States v. Guest}, 383 U.S. 745 (1966)}). Because state officials act under color of state law, the state-action requirement of \textit{Cruikshank} is satisfied, bringing both officials and private co-conspirators within federal enforcement power.
    \item \textbf{Fence-Sitter (Borderline / Contested):} \textit{Systematic State Failure to Prosecute Private Racial Violence.} Suppose state prosecutors and juries systematically refuse to indict or convict white mobs who terrorize Black citizens, effectively leaving Black residents without protection of state criminal law, though no state official directly joins the attacks. Under a strict reading of \textit{Cruikshank}'s state-action language, private violence remains outside federal reach; yet as historians and scholars note, \textit{Cruikshank} left unresolved whether a state's systematic \textit{omission} or refusal to afford equal protection itself constitutes ``state action'' justifying federal intervention under Section~5 of the Fourteenth Amendment (\href{https://civics.supremecourthistory.org/article/united-states-v-cruikshank/}{Supreme Court Historical Society}).
\end{enumerate}

\section{Critique}

\textbf{Originalist and Textualist Critiques:} Modern originalist scholars---and Supreme Court Justices---heavily criticize \textit{Cruikshank} for extending the \textit{Slaughter-House Cases}' evisceration of the Fourteenth Amendment's Privileges or Immunities Clause. Originalists such as Randy Barnett, Evan Bernick, and Ilya Shapiro, echoing Justice Clarence Thomas's concurrence in \href{https://www.law.cornell.edu/supct/html/08-1521.ZO.html}{\textit{McDonald v. City of Chicago}, 561 U.S. 742 (2010)}, argue that the Reconstruction Framers understood the Privileges or Immunities Clause to secure fundamental individual rights---explicitly including the right to keep and bear arms for self-defense---against state abridgment (\href{https://www.cato.org/sites/cato.org/files/2020-04/shapiro-blackman-gmlr-priv-imm-clause.pdf}{Shapiro \& Blackman, Cato Institute}; \href{https://illinoislawreview.org/online/what-cruikshank-really-means-for-the-second-amendment/}{Illinois L.~Rev.~Online}). By mechanically citing pre-Civil War precedent (\textit{Barron v. Baltimore}) as though the Fourteenth Amendment had changed nothing about the Bill of Rights, Waite's opinion broke sharply with the Amendment's original public meaning. At the same time, some textualist scholars defend \textit{Cruikshank}'s state-action holding on the ground that Section~1's text expressly begins with ``No State shall,'' distinguishing state denial of rights from private crime (\href{https://law.hofstra.edu/pdf/academics/journals/lawreview/lrv_issues_v34n04_cc2.huhn.final.pdf}{Hofstra L.~Rev.}).

\textbf{Progressive, Historical, and Doctrinal Critiques:} Historians and progressive constitutional scholars view \textit{Cruikshank} as a devastating judicial betrayal of Reconstruction. James Gray Pope argues that \textit{Cruikshank} ``belongs at the heart of the American constitutional canon'' because it effectively immunized white supremacist paramilitary terror against Black citizens, dismantled federal enforcement of the Fourteenth and Fifteenth Amendments, and paved the way for Jim Crow disenfranchisement (\href{https://journals.law.harvard.edu/crcl/wp-content/uploads/sites/80/2011/09/385_Pope.pdf}{Pope, Harv.~C.R.-C.L.~L.~Rev.}). Critics further point out the glaring formalism of Waite's remark that the Court could only ``suspect'' race motivated the Colfax Massacre---an atrocity universally understood as a racially targeted political slaughter---and note that the Court ignored the original understanding that a state's failure to protect Black citizens from organized terror was itself a denial of ``the equal protection of the laws'' (\href{https://www.cambridge.org/core/journals/law-and-society-review/article/judicial-abandonment-of-blacks-rethinking-the-state-action-cases-of-the-waite-court/6EFE4577A05FAA2791AB1541BFBFC371}{Law \& Society Review}; \href{https://www.fjc.gov/history/exhibits/circuit-court-opinions/United-States-v-Cruikshank}{Federal Judicial Center}).

\textbf{Underlying Competing Values:} \textit{Cruikshank} embodies the stark post-Civil War clash between \textbf{traditional state-centered federalism} (preserving states' primary police power and preventing Congress from converting all local crimes into federal offenses) and \textbf{national enforcement of equal citizenship} (vindicating the Reconstruction Amendments' promise that the federal government would protect freedpeople's civil and political rights when states could not or would not do so).

\section{Key Quotations}

\begin{quote}
``The same person may be at the same time a citizen of the United States and a citizen of a State, but his rights of citizenship under one of those governments will be different from those he has under the other.''
\end{quote}
\noindent --- Chief Justice Morrison R. Waite, Opinion of the Court (\href{https://supreme.justia.com/cases/federal/us/92/542/}{92 U.S. at 549}; \href{https://www.courtlistener.com/opinion/89309/united-states-v-cruikshank/}{CourtListener}).

\begin{quote}
``This is not a right granted by the Constitution. Neither is it in any manner dependent upon that instrument for its existence. The second amendment declares that it shall not be infringed, but this, as has been seen, means no more than that it shall not be infringed by Congress.''
\end{quote}
\noindent --- Chief Justice Morrison R. Waite, Opinion of the Court (\href{https://supreme.justia.com/cases/federal/us/92/542/}{92 U.S. at 553}).

\begin{quote}
``The Fourteenth Amendment prohibits a State from depriving any person of life, liberty, or property without due process of law, but this adds nothing to the rights of one citizen as against another. It simply furnishes an additional guaranty against any encroachment by the States upon the fundamental rights which belong to every citizen as a member of society.''
\end{quote}
\noindent --- Chief Justice Morrison R. Waite, Opinion of the Court (\href{https://supreme.justia.com/cases/federal/us/92/542/}{92 U.S. at 554}).

\begin{quote}
``We may suspect that race was the cause of the hostility, but it is not so averred. This is material to a description of the substance of the offence, and cannot be supplied by implication. Everything essential must be charged positively, and not inferentially.''
\end{quote}
\noindent --- Chief Justice Morrison R. Waite, Opinion of the Court (\href{https://supreme.justia.com/cases/federal/us/92/542/}{92 U.S. at 556}).

\begin{quote}
``I concur that the judgment in this case should be arrested, but for reasons quite different from those given by the court. \dots{} Certain other causes for arresting the judgment are assigned in the record which deny the constitutionality of the Enforcement Act; but, having come to the conclusion that the indictment is insufficient, it is not necessary to consider that question.''
\end{quote}
\noindent --- Justice Nathan Clifford, Dissenting (\href{https://supreme.justia.com/cases/federal/us/92/542/}{92 U.S. at 559, 569}).

\end{document}
